% Problem record op_95fccb9df34a08b1. This TeX file is derived from
% database/problems_json/op_95fccb9df34a08b1.json by scripts/sync-tex.mjs; edit the JSON record
% and rerun the script rather than editing this file.
\section{Parity is not in QAC0}
\label{sec:op_95fccb9df34a08b1}

\paragraph{Problem.}
Can polynomial-size constant-depth $\mathsf{QAC}^0$ circuits compute the
parity function, or equivalently, is
$\mathrm{PARITY}\notin\mathsf{QAC}^0$?

For $x=(x_1,\ldots,x_n)\in\{0,1\}^n$, define
\begin{equation}
  \mathrm{PARITY}_n(x)=x_1\oplus x_2\oplus\cdots\oplus x_n.
  \label{eq:95fc-parity}
\end{equation}
A $\mathsf{QAC}^0$ circuit family consists of polynomial-size quantum
circuits of constant depth built from arbitrary single-qubit gates and
generalized Toffoli gates, with polynomially many ancilla qubits initialized
to a fixed computational-basis state.

The conjecture is that there do not exist a constant $d$, a polynomial $p$,
and a family of depth-at-most-$d$ $\mathsf{QAC}^0$ circuits $C_n$ of size
and ancilla count at most $p(n)$ such that, for every $n$ and every
$x\in\{0,1\}^n$, measurement of a designated output qubit gives
\begin{equation}
  \Pr\bigl[\,C_n(x)\text{ outputs }\mathrm{PARITY}_n(x)\,\bigr]=1.
  \label{eq:95fc-compute}
\end{equation}
Equivalently, the conjecture asserts that no polynomial-size
constant-depth $\mathsf{QAC}^0$ family satisfies
\eqref{eq:95fc-compute} for the function defined in
\eqref{eq:95fc-parity}.

\subsection*{Status}

Unsolved

\subsection*{Source}

Moore explicitly posed the question whether ordinary $\mathsf{QAC}^0$ can
implement fanout; his constant-depth equivalence between fanout and parity
makes this the foundational form of the problem
\sourcecite{ref:95fc-moore}{Moo99}.  Later literature explicitly formulates
the expected negative answer as the conjecture that
$\mathrm{PARITY}\notin\mathsf{QAC}^0$, and the statement above is a
contributor-normalized formulation of that conjecture.

\subsection*{Progress}

\begin{itemize}
  \item Moore introduced $\mathsf{QAC}^0$ as a quantum analogue of classical
$\mathsf{AC}^0$ and showed that constant-depth quantum fanout and parity are
equivalent up to constant-depth reductions; determining whether fanout can be
implemented in ordinary $\mathsf{QAC}^0$ is equivalent to determining whether
parity belongs to $\mathsf{QAC}^0$
\sourcecite{ref:95fc-moore}{Moo99}.

  \item Rosenthal established strong lower bounds for depth-two $\mathsf{QAC}^0$
circuits approximating parity and gave the principal known constant-depth
upper bound for approximating parity with superpolynomially many resources;
polynomial-size arbitrary-constant-depth circuits remain unresolved
\sourcecite{ref:95fc-rosenthal}{Ros21}.

  \item Nadimpalli, Parham, Vasconcelos, and Yuen defined the Pauli spectrum of a
$\mathsf{QAC}^0$ circuit and proved a quantum analogue of low-degree Fourier
concentration for classical $\mathsf{AC}^0$, together with strong parity
lower bounds, when the number of auxiliary qubits is sufficiently restricted
as a function of circuit depth
\sourcecite{ref:95fc-npv}{NPVY24}.

  \item For depth $d$, Anshu, Dong, Ou, and Yao proved that computing functions of
linear approximate degree, including parity, requires at least
$n^{1+3^{-d}}$ ancilla qubits; since $\mathsf{QAC}^0$ permits an arbitrary
polynomial number of ancillae, this does not by itself resolve the
conjecture, but a sufficiently stronger depth-dependent exponent would imply
$\mathrm{PARITY}\notin\mathsf{QAC}^0$
\sourcecite{ref:95fc-adoy}{ADOY25}.

  \item Joshi, Tal, Vasconcelos, and Wright proved that depth-three $\mathsf{QAC}^0$
circuits cannot compute parity, regardless of their size or number of ancilla
qubits, and proved stronger approximation lower bounds at depth two; the
separation is therefore established for depth at most three
\sourcecite{ref:95fc-jtvw}{JTVW26}.

  \item Gretta, Gupta, and Joshi proved that if a polynomial-size $\mathsf{QAC}^0$
circuit family has nonzero Fourier mass at sufficiently high levels, then
parity can be computed exactly in $\mathsf{QAC}^0$; in their formulation the
corresponding Fourier-concentration statement is equivalent to separating
parity from $\mathsf{QAC}^0$
\sourcecite{ref:95fc-ggj}{GGJ26}.

  \item Sufficiently noticeable correlation with parity can be amplified and, with
constant-factor depth overhead and polynomial resources in the relevant
regime, converted into exact parity computation; average-case or
approximation lower bounds can therefore contribute directly to resolving
the exact conjecture
\sourcecite{ref:95fc-jtvw}{JTVW26}.
\end{itemize}

\subsection*{References}

\begin{enumerate}[label={},leftmargin=4.8em,labelsep=0.5em]
  \footnotesize
  \item[\textup{[Moo99]}]\label{ref:95fc-moore}
  Cristopher Moore, ``Quantum Circuits: Fanout, Parity, and Counting''
(1999).
\href{https://arxiv.org/abs/quant-ph/9903046}{arXiv:quant-ph/9903046}.

  \item[\textup{[Ros21]}]\label{ref:95fc-rosenthal}
  Gregory Rosenthal, ``Bounds on the QAC$^0$ Complexity of Approximating
Parity,'' in \emph{12th Innovations in Theoretical Computer Science
Conference (ITCS 2021)}, LIPIcs \textbf{185}, 32:1--32:20 (2021).
\href{https://doi.org/10.4230/LIPIcs.ITCS.2021.32}{doi:10.4230/LIPIcs.ITCS.2021.32};
\href{https://arxiv.org/abs/2008.07470}{arXiv:2008.07470}.

  \item[\textup{[NPVY24]}]\label{ref:95fc-npv}
  Shivam Nadimpalli, Natalie Parham, Francisca Vasconcelos, and Henry Yuen,
``On the Pauli Spectrum of QAC0,'' in \emph{Proceedings of the 56th Annual
ACM Symposium on Theory of Computing (STOC 2024)}, 1498--1506 (2024).
\href{https://doi.org/10.1145/3618260.3649662}{doi:10.1145/3618260.3649662};
\href{https://arxiv.org/abs/2311.09631}{arXiv:2311.09631}.

  \item[\textup{[ADOY25]}]\label{ref:95fc-adoy}
  Anurag Anshu, Yangjing Dong, Fengning Ou, and Penghui Yao, ``On the
Computational Power of QAC0 with Barely Superlinear Ancillae,'' in
\emph{Proceedings of the 57th Annual ACM Symposium on Theory of Computing
(STOC 2025)}, 1476--1487 (2025).
\href{https://doi.org/10.1145/3717823.3718189}{doi:10.1145/3717823.3718189};
\href{https://arxiv.org/abs/2410.06499}{arXiv:2410.06499}.

  \item[\textup{[JTVW26]}]\label{ref:95fc-jtvw}
  Malvika Raj Joshi, Avishay Tal, Francisca Vasconcelos, and John Wright,
``Improved Lower Bounds for QAC0,'' in \emph{Proceedings of the 58th Annual
ACM Symposium on Theory of Computing (STOC 2026)}, 2199--2209 (2026).
\href{https://doi.org/10.1145/3798129.3800922}{doi:10.1145/3798129.3800922};
\href{https://arxiv.org/abs/2512.14643}{arXiv:2512.14643}.

  \item[\textup{[GGJ26]}]\label{ref:95fc-ggj}
  Lucas Gretta, Meghal Gupta, and Malvika Raj Joshi, ``Parity
$\notin$ QAC0 $\Longleftrightarrow$ QAC0 is Fourier-Concentrated,'' in
\emph{67th Annual IEEE Symposium on Foundations of Computer Science (FOCS
2026)} (2026).
\href{https://arxiv.org/abs/2604.02793}{arXiv:2604.02793}.
\end{enumerate}

\subsection*{Comment}

The conjecture remains open for every sufficiently large constant depth
and in the regime of arbitrary polynomial ancilla count.  Known techniques
separate parity from $\mathsf{QAC}^0$ only in restricted regimes: depth at
most three, bounded ancilla, or approximation.  In particular, the
depth-dependent ancilla lower bound (Progress above) does not rule out
polynomial ancillae, and no Fourier-concentration lower bound is known for
$\mathsf{QAC}^0$ circuits with unrestricted ancilla.  The equivalence of
Gretta, Gupta, and Joshi (Progress above) shows that the conjecture is not
merely one algorithmic route but a precisely equivalent open question.

\subsection*{Field}

Quantum algorithm

\subsection*{Topic}

Quantum circuit complexity

\subsection*{ID}

\texttt{op\_95fccb9df34a08b1}
